% =====================================================================
% paper/sections/frontier_synthesis_group_size.tex
%
% PROVENANCE. This section distills an EXTERNAL frontier-model
% cross-examination of Pillar 3 (group-size effect and the
% GRPO<->DPO equivalence). The material below comes from a live
% 30-round adversarial cross-examination of ChatGPT Pro Extended and
% Gemini Deep Think (P3 rounds 3, 8, 9, 17, 23, 27 plus shared
% cross-cutting rounds 13, 14, 19, 20, 29, 30). These are REASONING
% CONTRIBUTIONS -- proposed theorems, closed forms, and falsifiable
% tests produced by the frontier models -- NOT experiments this paper
% ran. Everything is attributed as external synthesis and, where it
% reinforces or challenges our own measured G-sweep and "GRPO is
% secretly DPO" results (Section~\ref{sec:group-size}), the connection
% is stated in prose.
% Source digest: frontier_calls/digests/frontier_P3.md
% =====================================================================

\subsection{Frontier-Model Cross-Examination: Group Size and GRPO$\leftrightarrow$DPO}
\label{sec:frontier-group-size}

To stress-test the group-size analysis beyond our own measurements, we
submitted the Pillar~3 construction to an external cross-examination of
two frontier reasoning systems (ChatGPT Pro Extended and Gemini Deep
Think). The exchange is a source of \emph{proposed theorems and
falsifiable tests}, not additional experiments; we report the sharpest
contributions and, where they bear on our evidence---the measured
$100.3\%$ retention of $G{=}2$ vs $G{=}16$ on the near-ceiling
arithmetic task, the reconstructed $0.24$ absolute $G{=}4\!\to\!G{=}32$ swing at
canonical scale, and the sub-$\sqrt{G}$ SNR
(Table~\ref{tab:groupsize-iter27-synthesis})---we say so explicitly.

\paragraph{A proposed contrast-projection theorem: GRPO as an endogenous all-pairs DPO (frontier synthesis).}
Writing $s_i = \nabla_\theta \log \pi_\theta(y_i \mid x)$ for the
per-rollout score and $r_i$ for the reward, the frontier synthesis
derives the exact identity
\begin{equation}
  \sum_{i}(r_i - \bar r)\, s_i \;=\; \frac{1}{G}\sum_{i,j}(r_i - r_j)\, s_i,
  \label{eq:fs-contrast}
\end{equation}
so the mean-centred GRPO reward-contrast term has an all-pairs preference
(DPO) form in which every ordered pair contributes its reward margin
$r_i-r_j$. Two consequences follow. First, any group with zero reward
variance ($r_i \equiv r_j$) is projected out of that centered term; this
does not assert a zero total gradient. This is the algebraic origin of the zero-variance floor
$\mathrm{ZVF}(p,G)=p^G+(1-p)^G$ used throughout
Section~\ref{sec:group-size}. Second, unlike DPO, GRPO supplies its own
margins: it self-induces a preference signal \emph{only} on
mixed-reward groups. For binary rewards with $K$ successes the frontier
models give the closed form
$g_\text{GRPO}(x) = \tfrac{1}{G\,K(G-K)}\, g_\text{pair}(x)$, exposing
the $K(G-K)$ pair count as the mechanistic driver. In an external
cross-examination the models argued this is precisely why our measured
$G{=}2$ run retains $100.3\%$ of $G{=}16$ on the near-ceiling task: at
$p\!\to\!1$ almost every non-degenerate group already yields one usable
win--loss pair, so extra rollouts buy only redundant pairs.

\paragraph{Whitening turns the baseline into an adaptive-margin auto-curriculum.}
Gemini Deep Think separates GRPO's centering from its $1/\sigma$
whitening. Centering alone is a scaled leave-one-out (LOO) control
variate, $r_i - \bar r = \tfrac{G-1}{G}\,(r_i - b_{-i})$ with
$b_{-i}=\tfrac{1}{G-1}\sum_{j\ne i} r_j$, which is unbiased but injects
an $O(1/G)$ baseline-estimation penalty
$\mathrm{Var}(b_{-i}) = p_x(1-p_x)/(G-1)$. Dividing by the reward
standard deviation then \emph{breaks} control-variate neutrality: for a
group with $K$ successes the whitened advantages collapse to the
deterministic DPO-style scalars $A_+ = \sqrt{(G-K)/K}$ and
$A_- = -\sqrt{K/(G-K)}$, with adaptive margin
\begin{equation}
  \Delta_\text{margin} = A_+ - A_- = \frac{G}{\sqrt{K(G-K)}}.
  \label{eq:fs-margin}
\end{equation}
The frontier reading is that GRPO is a self-scheduling curriculum:
balanced groups ($K{=}G/2 \Rightarrow \Delta{=}2$) are gently throttled
while rare breakthroughs ($K{=}1 \Rightarrow \Delta \approx \sqrt{G}$)
receive amplified gradient. This rationalizes our measured
sub-$\sqrt{G}$ SNR ($1.46 \to 2.16$ from $G{=}2$ to $G{=}16$, only
$52\%$ of the $\sqrt{G}$ law, Table~\ref{tab:groupsize-iter27-synthesis}): the
$O(1/G)$ LOO penalty and the whitening non-linearity together cap
variance reduction well below the i.i.d.\ $\sqrt{G}$ ideal (frontier
synthesis).

\paragraph{Signal yield $Y_G(p)$ and why ``bigger $G$'' is not monotone.}
The frontier models formalize the useful-signal rate
\begin{equation}
  Y_G(p) = 1 - p^G - (1-p)^G,
  \label{eq:fs-yield}
\end{equation}
which is \emph{identical} to our gradient-utilization closed form
$\mathrm{GU}(p,G) = 1 - \mathrm{ZVF}(p,G)$: their ``contrastive yield''
and our ``gradient utilization'' are the same object reached from the
preference and the variance side respectively. Because $Y_G$ saturates
at $1$, increasing $G$ trades three quantities against one another---pair
generation ($\uparrow$), baseline variance ($\downarrow$ as $1/G$), and
inference compute ($\uparrow$ linearly)---so the compute-optimal group
size is finite, $G^\star = \arg\max_G Y_G(p)/G$, not ``as large as
possible''. For prompts bounded away from the extremes the frontier
bound $G \gtrsim \log\epsilon / \log(1-p_\text{min})$ (e.g.\
$p_\text{min}{=}0.1 \Rightarrow G\approx 29$ for $\mathrm{ZVF}<0.05$)
predicts that $G{\approx}32$ is the first genuinely DPO-equivalent
regime while $G{=}2$ is pairwise-equivalent only on already-unlocked
prompts. This reinforces our budget-dependent optimum
(Table~\ref{tab:groupsize-summary}) and gives a first-principles account
of why the Wu et al.\ $G{=}2{\sim}G{=}16$ equivalence fails to
generalize to $G{=}4{\sim}G{=}32$ at canonical scale
(Table~\ref{tab:groupsize-g4-g32}).

\paragraph{A decisive test separating variance-reduction from ZVF-reduction.}
Our data show a large held-out swing with $G$ but cannot, on their own,
say \emph{why}. The frontier synthesis proposes a $2{\times}2$
super-group design that separates the two candidate mechanisms while
holding tokens, steps, and KL fixed: Arm~A (natural small $G$), Arm~B
(variance-only: the same set of unlocked prompts but all $K(G{-}K)$
pairs used), Arm~C (ZVF-only: more unlocked prompts but a single random
win--loss pair each, or noise injected to match Arm~A's gradient
covariance), and Arm~D (natural large $G$). With a target swing of
$\Delta_{32-2}\approx 0.020$, the preregistered rule is
\emph{ZVF-dominant} if Arm~C recovers $\ge 75\%$ of the gain while
Arm~B recovers $\le 25\%$, and variance-dominant if the pattern
reverses. A complementary single-batch \emph{residual-isomorphism}
check subtracts a synthetic margin-scaled DPO gradient
$V_\text{mDPO}=\tfrac{1}{G}\sum_{i<j}(r_i-r_j)(s_i-s_j)$ from the
empirical GRPO gradient and tests whether the residual is pure
KL-penalty: $\cos(V_\text{GRPO}-V_\text{mDPO},\,V_\text{KL})\approx 1$
would confirm Eq.~\eqref{eq:fs-contrast} exactly, whereas $<0.99$ would
prove the group baseline carries stabilizing structure orthogonal to
pairwise preference. On the mechanism, the frontier models lean toward
ZVF-reduction: Gemini argued the $G{=}4\!\to\!G{=}32$ retention gap is
concentrated in high-accuracy bins ($+0.253$ at $p\ge 0.75$ versus
$+0.010$ at $p<0.5$), which is the signature of escaping the ZVF floor
rather than of classical $\sqrt{G}$ variance reduction. This
\emph{challenges} the variance-reduction reading we offer for the
$G{=}32$ reconstruction in Section~\ref{sec:group-size} and is a concrete,
falsifiable next experiment.

\paragraph{The Baseline-Equivalence-Index (BEI).}
Finally, the frontier synthesis proposes a Baseline-Equivalence-Index
to pin \emph{when} the GRPO group mean, a PPO value head, and a DPO
margin are interchangeable. Rather than compare baseline mean-squared
errors it compares the resulting actor \emph{updates} $g_G$ (group
baseline) and $g_V$ (value head),
\begin{equation}
  \mathrm{BEI}_G =
  \frac{E_x\langle g_V, g_G\rangle}
       {\sqrt{E_x\|g_V\|^2\; E_x\|g_G\|^2}}
  \cdot
  \min\!\Big(\tfrac{E\|g_G\|}{E\|g_V\|},\,\tfrac{E\|g_V\|}{E\|g_G\|}\Big),
  \label{eq:fs-bei}
\end{equation}
so that $\mathrm{BEI}\approx 1$ requires the two updates to match in
\emph{both} direction and scale---a value head can match direction
while silently rescaling the effective learning rate, which a cosine
alone would miss. The proposed boundary is exchangeable above
$\approx 0.95$, distinct estimators below $\approx 0.8$ or on a
projection sign flip; equivalently, the group mean overtakes a trained
critic once its $O(1/G)$ Monte-Carlo variance drops below the critic's
mean-squared error. Because our measured SNR rises only sublinearly
($52\%$ of $\sqrt{G}$), the frontier prediction is that BEI should climb
sublinearly and \emph{saturate before} critic-free GRPO becomes fully
value-head-equivalent at the group sizes we test---i.e.\ $G{=}16$ is not
yet in the PPO-equivalent regime. At the opposite end the same account
recovers the exact GRPO$\leftrightarrow$DPO bridge: at $G{=}2$ the
centred update assigns $+\tfrac12$ to the winner and $-\tfrac12$ to the
loser, precisely a one-pair DPO step.

\paragraph{Bottom line (frontier synthesis).}
Taken together, the external analysis recasts the group-size knob as a
preference-learning dial: $G$ controls how many endogenous DPO pairs a
prompt contributes (Eq.~\ref{eq:fs-contrast}), the whitened margin turns
that count into an auto-curriculum (Eq.~\ref{eq:fs-margin}), and the
useful fraction $Y_G(p)=\mathrm{GU}(p,G)$ (Eq.~\ref{eq:fs-yield}) bounds
the payoff so that the compute-optimal $G$ is finite. This is consistent
with our measured near-ceiling equivalence at small $G$ and with the
failure of that equivalence at canonical scale, while sharpening the one
question our data leave open---whether the large-$G$ win is
variance-reduction or ZVF-escape---into a preregistered, falsifiable
test. We stress again that these are external reasoning contributions,
not experiments this paper ran.
